% !TeX root = ../../金融经济学_王江_讲义及习题.tex
% 知识点：references/金融经济学/金融经济学-王江_261013.pdf，印刷页 87–95
% 课堂补充：references/金融经济学/金融经济学_第6章 期望效用函数.pptx（夏晖，2026 年 10 月），共 34 页；含图片公式核对。
% 习题：references/金融经济学/史上最全最完整的金融经济学答案(王江版).pdf，页 26–27（印刷页 25–26）

\lecture{期望效用函数}

\ifdefined\WJHandoutKnowledge

无套利给出的是证券价格之间的关系：由可交易证券的价格和支付可以定出状态价格的集合，也可以给冗余证券定价。它几乎不依赖经济的具体特征，因此也回答不了这些价格本身由什么决定。以二叉树为例，无套利能给出期权价格以及 $u>1+r_F>d$，但不能告诉股票过程和无风险利率从何而来，也不能把所有证券的价格都定出来。要追问价格水平以及证券在资源配置中的作用，必须回到第 2 章的均衡分析，而对偏好、风险分布和市场结构作出更具体的假设。本章从偏好开始。

第 2 章用不满足性、连续性和凸性等条件刻画偏好；一般效用表示的拟凹性与本章伯努利效用的凹性需要区分，$U(c)$ 的具体形式仍很宽。一般的 $U(c)$ 同时依赖两件事：未来状态的概率，以及各状态下对消费的偏好。本章用\textbf{独立性公理}把二者分开，得到期望效用；再加上状态独立、时间可加等假设，得到后面各章使用的形式；最后指出这些假设排除了哪些行为，以及文献里常见的拓展。

\begin{tcolorbox}[knowledge,title=6.1 期望效用函数]
先看 1 期没有不确定性、两期消费都确定的情形，效用记为 $u(c_0,c_1)$。若 1 期有两个等可能状态 $a$ 与 $b$，消费分别为 $c_{1a}$、$c_{1b}$。状态 $a$ 下消费路径 $(c_0,c_{1a})$ 的效用记为 $u_a(c_0,c_{1a})$，下标表示这条路径的效用可以依赖状态。一个直接的做法是用概率作权数，取各消费路径效用的平均值：
\[
    \tfrac12 u_a(c_0,c_{1a})+\tfrac12 u_b(c_0,c_{1b}).
\]
一般地，若上述分离成立，效用就是各消费路径效用的期望，
\[
    U(c)=\sum_{\omega\in\Omega}p_{\omega}\,u_{\omega}(c_0,c_{1\omega}).
    \tag{6.1}
\]
教材 (6.1) 把概率记为 $\pi_\omega$。这种 $U$ 称为\textbf{期望效用函数}（expected utility function）。概率以线性权重进入总效用；本书另假定这些概率外生给定、为所有人共同知晓，对消费本身的喜好则因人而异。期望效用也可配合不同人的主观概率使用，“共同概率”并不是期望效用形式本身强制的。它还蕴涵下面的公理。

\begin{tcolorbox}[knowledge]
\textbf{公理 4（独立性公理）}\quad 设消费计划 $c$ 与 $c'$ 在某一状态上的消费路径同为 $x$，且 $c\succeq c'$。把这个共同路径 $x$ 换成另一路径 $y$ 后，$c$ 与 $c'$ 的排序不变。
\end{tcolorbox}

含义：若 $c$ 与 $c'$ 只在状态 $\omega$ 上同为路径 $x$，且 $c\succeq c'$，则剩余状态下路径的综合比较支持 $c$，不要求它在每一个剩余状态都更好。把 $x$ 换成 $y$ 后排序不变，说明其他状态下的比较不受状态 $\omega$ 上发生什么的影响。各互斥状态下的效用彼此独立。

两状态、等概率时可以直接验证。设
\[
    c=\begin{bmatrix}c_0\\ c_{1a}\\ c_{1b}\end{bmatrix},\qquad
    c'=\begin{bmatrix}c_0\\ c'_{1a}\\ c_{1b}\end{bmatrix},
\]
二者在状态 $b$ 上的路径相同。若 $c\succeq c'$，则
\[
    U(c)=\tfrac12\bigl[u_a(c_0,c_{1a})+u_b(c_0,c_{1b})\bigr]
    \ge\tfrac12\bigl[u_a(c_0,c'_{1a})+u_b(c_0,c_{1b})\bigr]=U(c').
\]
把状态 $b$ 的共同路径换成 $(c_0,c'_{1b})$，得到 $c''$、$c'''$。$u_b$ 在两边的贡献仍相同，故 $c''\succeq c'''$。

各消费路径与互斥的状态相对应：\textbf{一个状态发生意味着其余状态不发生，因此不同路径之间可以分开比较}。这是独立性公理说得通的一个理由。

% PPT 第 6–8 页：替换共同后果的含义与数字验证。
\textbf{独立性的对象是偏好排序。}这里将两个计划在某状态的\textbf{共同后果同时替换}，其余不同部分保持不变；它不表示“只改变一个计划的消费也不影响排序”，也不要求证券收益等随机变量在统计上独立。互斥性说明为何可以考虑这种分离，但独立性本身仍是一项可检验、也可能被违反的偏好假设。

\textbf{课堂数字验证（PPT 第 8 页）：}固定相同的当前消费，未来两状态等概率，$u(x)=\sqrt{x}$。原来的未来消费为 $c=(100,16)$、$c'=(64,16)$；把共同的雨天消费 16 同时换成 36，得到 $c''=(100,36)$、$c'''=(64,36)$。
\[
\begin{array}{c|cc|c}
 \text{计划}&\text{晴天消费}&\text{雨天消费}&\text{未来期望效用}\\\hline
 c&100&16&\tfrac12(10+4)=7\\
 c'&64&16&\tfrac12(8+4)=6\\
 c''&100&36&\tfrac12(10+6)=8\\
 c'''&64&36&\tfrac12(8+6)=7
\end{array}
\]
两个比较的效用差均为 $\tfrac12(10-8)=1$。共同部分改变了总效用水平，却没有改变两计划之间的排序。


\begin{tcolorbox}[example]
\begin{question}
\textbf{例 6.1}\quad 明天或晴或雨。计划 $c$：晴则在海滩玩 $4$ 小时，雨则看 $4$ 小时电视。计划 $c'$：晴则海滩 $2$ 小时再看电视 $2$ 小时，雨则仍看 $4$ 小时电视。这里出现两种消费品（海滩与电视），但比较的仍是整条消费计划。参与者更喜欢海滩，故晴天时更喜欢 $c$ 而不是 $c'$，即 $c\succeq c'$。

另给两个计划：$\tilde{c}$ 为晴则海滩 $4$ 小时、雨则海滩 $4$ 小时；$\tilde{c}'$ 为晴则海滩 $2$ 小时再电视 $2$ 小时、雨则海滩 $4$ 小时。独立性公理要求他选择 $\tilde{c}$ 而不是 $\tilde{c}'$。只要雨天的安排相同，雨天具体做什么不影响晴天的偏好。
\end{question}
\begin{solution}
两个计划在雨天的消费路径相同，排序已由晴天的路径决定。把雨天的共同路径从“看电视”换成“去海滩”，晴天的比较原样保留。
\end{solution}
\end{tcolorbox}

\begin{tcolorbox}[example]
\begin{question}
\textbf{例 6.2}\quad 仍在例 6.1 的晴天里比较两个消费组合，而不是两条跨状态的消费计划。组合 $c_1$：海滩 $4$ 小时并喝 $1$ 升水；组合 $c_2$：电视 $4$ 小时并喝 $1$ 升水。若他选择 $c_1$，把那 $1$ 升水去掉之后，是否仍选择海滩 $4$ 小时？
\end{question}
\begin{solution}
不一定。酷暑里他可能因此改选待在室内。海滩和水是互补品，总效用取决于两种消费品的搭配，而不只是单件消费品。独立性公理讲的是\textbf{互斥状态}之间的独立，不是同一状态里不同商品之间的独立。
\end{solution}
\end{tcolorbox}

独立性是期望效用形式的必要条件，也是充分条件。

\[
\text{独立性} \Leftrightarrow \text{期望效用}
\]

\begin{tcolorbox}[knowledge]
\textbf{定理 6.1（Debreu）}\quad 在独立性公理下，定理 2.1 给出的效用函数具有 (6.1) 的期望效用形式。
\end{tcolorbox}

解释：不确定状态下的效用，是每一可能状态下消费路径效用的加权平均，权重为该状态的概率。 教材在既定消费空间和前述偏好条件下陈述此表示结果；后文用到的彩票混合独立性还需以彩票为选择对象，不能把仅有连续性或互斥性本身当作期望效用表示的充分条件。
\end{tcolorbox}

\begin{tcolorbox}[knowledge,title=6.2 附加假设]
(6.1) 仍相当一般。下面几条简化是为了在最优消费、组合选择和均衡价格中得到更具体的结论。除特别声明外，后文都采用这些形式。

\tcbline

\textbf{A.\ 状态独立}

(6.1) 允许同一消费路径在不同状态下给出不同效用。这是否合理，取决于状态和消费怎样描述。用“晴天喝 $1$ 升水”和“多云喝 $1$ 升水”来写，效用除了水量还会依赖天气；若改用“阳光的消费量”刻画晴和阴，扣掉阳光之后，两个状态可以变成一样的，同时喝 $1$ 升水加享受 $1$ 小时阳光的效用就不必再依赖“晴或阴”这个标签。对市场总体而言状态仍有区别（可消费的阳光总量不同，会影响价格）；若只关心个人消费和效用，就可以认为没有区别。于是状态独立的期望效用是自然的：
\[
    u_{\omega}(c_0,c_{1\omega})=u(c_0,c_{1\omega}),
    \qquad
    U(c)=\sum_{\omega\in\Omega}p_{\omega}\,u(c_0,c_{1\omega}).
    \tag{6.2}
\]

\textbf{课堂解释（PPT 第 13 页）：}把完整消费篮子写为 $(x,s)$，$x$ 是冰淇淋数量，$s$ 是阳光数量。晴天可消费 $(1,1)$，阴天为 $(1,0)$；用同一个 $v(x,s)$，可以有 $v(1,1)=10$、$v(1,0)=5$。效用差异来自\textbf{消费篮子不同}，无需另给“晴天效用函数”和“阴天效用函数”。这不表示阳光可自由交易，也不意味着任何状态依赖都能在既定商品定义下消除；它说明状态独立性依赖模型怎样完整描述消费后果。

若不加限制，状态依赖的效用非常灵活，几乎任何行为都可以被解释成“他只是不喜欢在某个状态下消费”。这种写法方便，但假设与结论之间几乎没有距离，本书对此持审慎态度，后文默认 (6.2)。

\tcbline

\textbf{B.\ 时间可加性}

(6.2) 中，1 期的效用仍可以依赖 0 期消费，反之亦然，即交叉导数 $\partial^{2}u/(\partial c_0\partial c_{1\omega})$ 可以非零。实物消费里这很常见：今天吃了很多牛肉，明天对牛肉的胃口下降；今天多抽烟，明天更想抽。但本书配置的是财富和资源，而不是特定实物。为简单起见，假定
\[
    \frac{\partial^{2}u(c_0,c_1)}{\partial c_0\,\partial c_1}=0,
\]
从而
\[
    u(c_0,c_1)=u_0(c_0)+u_1(c_1).
\]
\textbf{为什么交叉偏导为零会得到可加形式？}在矩形消费定义域内，若 $u$ 二阶连续可微且 $\partial_{01}u=0$，则 $\partial_0u(c_0,c_1)$ 不随 $c_1$ 变化，可记为 $f(c_0)$。对 $c_0$ 积分，
\[
 u(c_0,c_1)=\int f(c_0)\,dc_0+g(c_1)=u_0(c_0)+u_1(c_1).
\]
反过来，可加形式的交叉偏导显然为零。它排除的是一期消费对另一期\textbf{边际效用}的影响，并不表示两期消费在预算中互不约束。

一条消费路径的效用等于各期效用之和，称为\textbf{时间可加}（time additive）或\textbf{时间可分}（time separable）。期望效用化为
\[
    U(c)=u_0(c_0)+\sum_{\omega\in\Omega}p_{\omega}\,u_1(c_{1\omega}).
    \tag{6.3}
\]

有时再设两期的效用函数只差一个正的倍数，
\[
    u_0(c_0)+u_1(c_1)=u(c_0)+\rho\,u(c_1),\qquad\rho>0.
\]
$\rho$ 称为\textbf{时间偏好系数}（time preference coefficient），也称为时间折扣系数。$\rho=1$ 时，参与者不关心消费发生在哪一期；$\rho<1$ 时，同样的消费越早越好，推迟一期使等量消费的效用降为原来的 $\rho$ 倍。此时
\[
    U(c)=u(c_0)+\rho\sum_{\omega\in\Omega}p_{\omega}\,u(c_{1\omega}).
    \tag{6.4}
\]
三个因素完全分开：状态概率 $P=\{p_\omega\}$、时间性 $\rho$、消费本身的效用 $u(\cdot)$。称呼上常把 $u(\cdot)$ 也叫效用函数，但须与总效用 $U$ 分开。

\textbf{与资产定价的联系（PPT 第 5、15 页）：}可加形式给出 $\partial_0U=u'(c_0)$、$\partial_\omega U=\rho p_\omega u'(c_{1\omega})$。内点最优时，状态价格满足
\[
 \phi_\omega=\rho p_\omega\frac{u'(c_{1\omega})}{u'(c_0)},\qquad
 m_\omega=\frac{\phi_\omega}{p_\omega}
 =\rho\frac{u'(c_{1\omega})}{u'(c_0)}.
\]
$m$ 是随机折现因子，证券定价写为 $S_n=\mathrm E[m\widetilde X_n]$。概率、时间偏好与消费边际效用的作用由此分别显现。若效用时间不可加，边际效用会依赖整条消费路径，推导仍可进行，但不再自动具有上述单期边际效用比的简洁形式；是否存在解析解取决于具体模型。


$U$ 是序数的：正单调变换不改变消费计划的排序，数值大小本身不重要。$u$ 不是这样。对 $u$ 作正单调变换，可能改变消费计划的排序。

\begin{tcolorbox}[example]
\begin{question}
\textbf{例 6.3}\quad 两个状态等可能。消费计划
\begin{center}
\begin{tabular}{@{}c@{\quad}c@{}}
\textbf{计划 A} & \textbf{计划 B} \\[0.35em]
\WJEventTree[left={$1$},stem]{1,1} &
\WJEventTree[left={$1$},stem]{3,0}
\end{tabular}
\end{center}
偏好由 $u(c_0)+\frac12\bigl[u(c_{11})+u(c_{12})\bigr]$ 表示。先取 $u(c)=c$，再取 $v(c)=\sqrt{u(c)}=\sqrt{c}$。两种效用函数下，A 与 B 的排序是否相同？
\end{question}
\begin{solution}
$u(c)=c$ 时，
\[
    U(A)=1+\tfrac12(1+1)=2,
    \qquad
    U(B)=1+\tfrac12(3+0)=\tfrac52.
\]
$B$ 优于 $A$。换成 $v(c)=\sqrt{c}$ 后，
\[
    U(A)=1+\tfrac12(\sqrt{1}+\sqrt{1})=2,
    \qquad
    U(B)=1+\tfrac12(\sqrt{3}+\sqrt{0})=1+\tfrac{\sqrt{3}}{2}.
\]
$1+\sqrt{3}/2<2$，故 $A$ 优于 $B$。排序反过来了。
\end{solution}
\end{tcolorbox}

\textbf{课堂对照（PPT 第 18 页）：}令 $C$ 确定获得 1.5，$D$ 等概率获得 1 或 2。在 $u(x)=x$ 下，两者效用都为 1.5；换成 $v(x)=\sqrt{x}$ 后，
\[
 V(C)=\sqrt{1.5}\approx1.224745,
 \qquad V(D)=\tfrac12(1+\sqrt2)\approx1.207107.
\]
$C$ 严格优于 $D$。虽然平方根保持每个确定结果的排序，却改变了对彩票的排序。对\textbf{总效用}变换成 $\sqrt{U}$ 与先对\textbf{单期效用}变换后再取期望不同：一般有
$\sqrt{\mathrm E[u(\widetilde W)]}\ne\mathrm E[\sqrt{u(\widetilde W)}]$。

因此 $u$ 具有\textbf{基数性}（cardinal）：不同消费水平上效用的相对差值会进入期望，从而影响排序。在期望效用形式下，偏好由 $u$ 唯一描述；$u$ 的正线性变换（$au+b$，$a>0$）与 $u$ 等价，一般的正单调变换则不行。这与定理 2.1 里序数的 $U$ 不同。 对 (6.4) 使用同一个变换 $v=au+b$，可直接验证
\[
 V(c)=aU(c)+b(1+\rho),\qquad a>0,
\]
所以所有计划的排序不变。“基数性”不表示效用具有可直接观察的绝对单位，也不保证不同人的效用数值可以直接比较。
\end{tcolorbox}

\begin{tcolorbox}[knowledge,title=6.3 期望效用函数的拓展]
(6.4) 加上单调和凹，对偏好的刻画简洁，也因此排除了不少行为。当由它推出的结论与观察不一致时，自然会放宽假设。\textbf{行为金融学}（behavioral finance）讨论这些实证偏离及其含义，超出本书范围。这里只说明它们怎样改写效用函数。例子包括（远不止这些）：

\begin{itemize}
    \item \textbf{习惯}（habit formation）：效用依赖消费相对于过去水平的增量，是时间不可分的一种形式。例如
    \[
        u(c_0,c_1)=u(c_0)+\rho\,u(c_1-hc_0),\qquad h\ge 0.
    \]
    1 期效用取决于相对于前期消费的增量。若采用要求正消费的对数形式，须有 $c_1-hc_0>0$，且接近该边界时效用趋于 $-\infty$。只有边际效用 $u'(0)=\infty$ 本身不能推出零消费效用是 $-\infty$，例如平方根效用在零处仍有限。
    \item \textbf{攀比}（catching up with the Joneses）：效用不仅依赖自己的消费，还依赖别人的消费，
    \[
        u(c_0,c_1)=u(c_0-kC_0)+\rho\,u(c_1-kC_1),\qquad 0\le k<1,
    \]
    其中 $C_0$、$C_1$ 为两期的人均总消费。实际效用取决于个人消费与人均消费之差。给定状态下总禀赋是给定的，这也可以看成状态依赖的一种特殊形式。
    \item \textbf{状态依赖}（state dependence）：例如效用直接依赖经济总体的消费水平。由第 3 章，边际效用高的状态会被配置更多资源。若“经济情况好”时边际效用更高，参与者会把更多资源放到好状态、给差状态留下更少，与通常把财富从好状态转移到坏状态的保险直觉相反。这类偏好可以解释为过于乐观，高估了好状态发生的可能，也是经济与行为金融讨论的对象。
    \item \textbf{一阶风险厌恶}（first order risk aversion）：效用在参考点处不可微。例如
    \[
        u(c)=
        \begin{cases}
            a_{+}(c-\bar{c}), & c\ge\bar{c},\\
            a_{-}(c-\bar{c}), & c<\bar{c},
        \end{cases}
        \qquad a_{-}>a_{+}>0.
    \]
    $u$ 在 $c=\bar{c}$ 处不可微，跨越参考点时损失与收益的边际价值不同。这可以刻画一种参考点不对称，但本例两侧均线性，并不是包含损失域凸性、概率加权等要素的完整前景理论。其一阶风险厌恶含义在下一章讨论。
    \item \textbf{不确定性厌恶}（uncertainty aversion）：可引出超出标准期望效用的模型。要区分已知概率下对风险的态度与概率本身不明确时的模糊厌恶；违反独立性公理不能自动等同于模糊厌恶。下节按 PPT 补充已知概率的彩票反例。
\end{itemize}

\textbf{习惯与攀比的区别（PPT 第 19 页）：}习惯的参照来自自己的过去，攀比的参照来自他人的消费。在习惯模型 $u(c_0)+\rho u(c_1-hc_0)$ 中，
\[
 \partial_0U=u'(c_0)-\rho h u'(c_1-hc_0),\qquad
 \partial_1U=\rho u'(c_1-hc_0).
\]
今天多消费不仅提供当期效用，还提高明天的习惯基准，因此边际收益要扣除后一项。攀比模型在个人把 $C_t$ 视为给定时，边际效用为 $u'(c_t-kC_t)$，基准来自社会环境。两类偏好可影响组合与价格，但并不单凭函数形式就必然推出“追涨杀跌”或“羊群行为”。

这些形式都还落在第 2 章的基本框架里，只是比 (6.3)、(6.4) 更灵活。本书此后只用 (6.3) 或 (6.4)。基本框架一旦清楚，再把效用换成上面的某一类拓展是直接的。
\end{tcolorbox}

\begin{tcolorbox}[knowledge,title=6.4 不确定性下的决策：圣彼得堡悖论]
% PPT 第 21–24 页，包含图片中的奖金表、级数与 CE。
掷一枚公平硬币，直到第一次正面出现。若首次正面在第 $i$ 次出现，奖金为 $X_i=2^{i-1}$，概率为 $2^{-i}$：
\[
\begin{array}{c|ccccc}
 \text{路径}&H&TH&TTH&TTTH&\cdots\\\hline
 \text{概率}&1/2&1/4&1/8&1/16&\cdots\\
 \text{奖金}&1&2&4&8&\cdots
\end{array}
\]
若只按期望奖金评价，
\[
 \mathrm E[X]=\sum_{i=1}^{\infty}2^{-i}2^{i-1}
 =\tfrac12\sum_{i=1}^{\infty}1=\infty.
\]
但无限的期望奖金并不意味着人们愿意支付无限入场费。课堂用凹效用说明，关心效用与只关心期望金额会给出不同评价。

\textbf{A.\ 对数效用的计算}\quad 暂把奖金本身视为最终财富，$u(x)=\ln x$。则
\[
 \mathrm E[u(X)]=\ln2\sum_{i=1}^{\infty}2^{-i}(i-1).
\]
从几何级数开始，在 $|x|<1$ 时逐项求导：
\[
 \sum_{i=0}^{\infty}x^i=\frac1{1-x},\qquad
 \sum_{i=1}^{\infty}ix^{i-1}=\frac1{(1-x)^2}.
\]
令 $x=1/2$ 并乘以 $1/2$，得到
\[
 \sum_{i=1}^{\infty}i2^{-i}=2,\qquad
 \sum_{i=1}^{\infty}2^{-i}=1.
\]
所以 $\mathrm E[\ln X]=\ln2(2-1)=\ln2$，确定性等值为
\[
 CE_{\ln}=\exp(\mathrm E[\ln X])=2.
\]

\textbf{B.\ 平方根效用的计算}\quad 若 $u(x)=\sqrt{x}$，
\[
\begin{aligned}
 \mathrm E[\sqrt X]
 &=\sum_{i=1}^{\infty}2^{-i}2^{(i-1)/2}
 =\frac12\sum_{j=0}^{\infty}(1/\sqrt2)^j\\
 &=\frac{1/2}{1-1/\sqrt2}=1+\frac1{\sqrt2}.
\end{aligned}
\]
故
\[
 CE_{\sqrt{\cdot}}=\left(1+\frac1{\sqrt2}\right)^2
 =\frac32+\sqrt2\approx2.914214.
\]
两种凹效用都使该彩票的期望效用和 CE 有限，但 CE 数值不同。凹性并不保证\textbf{任意}无界彩票的期望效用都有限；结论须针对具体分布与效用检验。

\textbf{CE 与入场费须区分：}上述计算没有纳入初始财富和支付费用。若初始财富为 $w$，入场费为 $f$，无差异条件应为
\[
 \sum_{i=1}^{\infty}2^{-i}u(w-f+2^{i-1})=u(w),
\]
且每个结果须在效用定义域内。因此不能把 2 或 2.914214 直接当作所有财富者的最高入场费。有限财富、主办者偿付能力与奖金上限也会影响实际决策。
\end{tcolorbox}

\begin{tcolorbox}[knowledge,title=6.4（续）阿莱斯悖论与独立性]
% PPT 第 25–27 页；p3/p4 沿用 PPT 第 26 页的编号。
奖金单位为\textbf{万美元}，四种彩票为：
\[
\begin{array}{c|ccc}
 &0&100&500\\\hline
 P_1&0&1&0\\
 P_2&0.01&0.89&0.10\\
 P_3&0.90&0&0.10\\
 P_4&0.89&0.11&0
\end{array}
\]
课堂讨论的一种选择模式是 $P_1\succ P_2$，同时 $P_3\succ P_4$。第一对偏好确定得到 100；第二对则偏好 10\% 得到 500，而非 11\% 得到 100。

\textbf{为什么不可能由同一个期望效用表示？}若 $P_1\succ P_2$，就有
\[
 u(100)>0.01u(0)+0.89u(100)+0.10u(500),
\]
即
\[
 0.11u(100)>0.01u(0)+0.10u(500).\tag{6.A}
\]
而 $P_3\succ P_4$ 要求
\[
 0.90u(0)+0.10u(500)>0.89u(0)+0.11u(100),
\]
即 $0.01u(0)+0.10u(500)>0.11u(100)$，与 (6.A) 矛盾。这与 $u$ 是否凹无关；仅改变风险厌恶程度不能消除这个矛盾。

\textbf{共同后果是怎样被替换的？}令 $Q$ 为“以 $1/11$ 概率得到 0、以 $10/11$ 概率得到 500”的彩票，$\delta_x$ 表示确定得到 $x$。则
\[
\begin{aligned}
 P_1&=0.11\delta_{100}+0.89\delta_{100},&
 P_2&=0.11Q+0.89\delta_{100},\\
 P_4&=0.11\delta_{100}+0.89\delta_0,&
 P_3&=0.11Q+0.89\delta_0.
\end{aligned}
\]
彩票混合独立性要求，对任意 $0<\alpha<1$，
\[
 L\succeq M\ \Longleftrightarrow\
 \alpha L+(1-\alpha)R\succeq\alpha M+(1-\alpha)R.
\]
把双方共有的 89\%“得到 100”同时换为“得到 0”，原来的排序应保持，故 $P_1\succ P_2$ 应推出 $P_4\succ P_3$。观察到的相反选择体现\textbf{共同后果效应或确定性效应}。这是已知概率的彩票，不是因为人们不知道概率才产生的悖论。

PPT 列出几种回应：坚持独立性作为规范要求；研究接近 0、1 的概率是否有特殊作用；引入后悔；或放宽独立性、构建非期望效用模型。悖论说明这组选择与特定公理不能同时成立，并不单凭一个例子就证明所有人、所有情形都遵循另一种模型。
\end{tcolorbox}

\begin{tcolorbox}[knowledge,title=6.4（续）秩依效用与前景理论]
\textbf{A.\ 秩依效用（PPT 第 28 页）}\quad 秩依效用保留最终结果的效用，但按结果排序后的\textbf{累计概率}构造决策权重。设 $x_1\le\cdots\le x_n$、对应概率为 $p_i$，取递增函数 $g:[0,1]\to[0,1]$，满足 $g(0)=0,g(1)=1$，可写成
\[
 V=\sum_{i=1}^n\pi_i u(x_i),\qquad
 \pi_i=g\!\left(\sum_{j=i}^np_j\right)
       -g\!\left(\sum_{j=i+1}^np_j\right).
\]
$\pi_i\ge0$ 且 $\sum\pi_i=1$；$g(p)=p$ 时退回期望效用。权重依赖结果的\textbf{秩}，不是把每个 $p_i$ 分别变换后直接相加。某些加权函数可刻画高估小概率等现象，但这些性质不是所有秩依效用模型都自动具备的；概率为零的事件仍有零权重。

\tcbline
\textbf{B.\ 前景理论（PPT 第 29–31 页）}\quad 前景理论的关键是\textbf{参考点、收益与损失的不对称及概率决策权重}。价值函数作用于相对参考点 $r$ 的变化 $x-r$，而不只作用于最终财富。常见示意形式为
\[
 v(y)=\begin{cases}y^\alpha,&y\ge0,\\-\lambda(-y)^\alpha,&y<0,\end{cases}
 \qquad0<\alpha<1,\quad\lambda>1.
\]
收益域凹、损失域凸；$\lambda>1$ 使同额损失的价值变化绝对值更大。损失厌恶与风险厌恶不是同一个概念：前者比较收益和损失，后者比较同一域内确定与不确定结果。完整模型还需指定概率权重；仅给这个价值函数并不能推出所有选择。

\textbf{课堂例：同样最终财富，不同参考点}\quad 明确两项选择是\textbf{最终财富}：$A$ 确定为 50 元，$B$ 等概率为 100 元或 0 元。设 $\alpha=1/2,\lambda=2$，暂用原概率权重，只考察参考点的影响。
\[
\begin{array}{c|cc|c}
 \text{参考点}&V(A)&V(B)&\text{排序}\\\hline
 r=0&\sqrt{50}\approx7.071&\tfrac12\sqrt{100}+\tfrac12\sqrt0=5&A\succ B\\
 r=100&-2\sqrt{50}\approx-14.142&\tfrac12v(0)+\tfrac12v(-100)=-10&B\succ A
\end{array}
\]
相同最终财富在第一种框架下是收益，在第二种框架下是损失，因而可能发生偏好反转。若题目改为“在现有财富之外\textbf{新增}50 元”，则必须先加回现有财富再减参考点，不能把新增收益和最终财富混用。
\end{tcolorbox}

\begin{tcolorbox}[example]
\textbf{课堂例：丢票与丢钱（PPT 第 32–33 页）}\quad 两张演出票总价为 $K$，仍能买到同样的票且没有流动性约束。比较：已经买好的票丢失后是否重买；与准备买票前丢失了 $K$ 元现金后是否仍买。两种情形的经济损失均为 $K$，再次购票都需另付 $K$。

如果只考虑最终财富、同样的演出体验和同样的后续机会，两种情形的可选结果相同，标准期望效用模型应给出相同的购票决定。但心理账户可能使参与者给出不同回答：丢票后的重买被计入“看演出”账户，感觉这场演出的成本变成 $2K$；丢现金则记入一般现金损失，演出账户仍只记 $K$，因而更容易接受购买。

这说明\textbf{损失怎样归类和描述}会影响选择。心理账户与参考点、框架效应相关，但不等于前景理论的全部定义；若两种情形实际存在不同的找回概率、票价或预算约束，也应先控制这些差异。
\end{tcolorbox}

\begin{tcolorbox}[knowledge,title=本章小结]
\begin{enumerate}
    \item 独立性公理下，效用可以写成各消费路径效用的期望 $U(c)=\sum_{\omega}p_{\omega}u_{\omega}(c_0,c_{1\omega})$。独立性是这一形式的充分必要条件（定理 6.1）。它针对的是互斥状态，不是同一状态中不同商品。
    \item 状态独立把 $u_{\omega}$ 收成同一个 $u$；时间可加再把它拆成 $u_0(c_0)+\sum p_{\omega}u_1(c_{1\omega})$。若两期只差时间偏好系数 $\rho$，则 $U(c)=u(c_0)+\rho\sum p_{\omega}u(c_{1\omega})$。总效用 $U$ 是序数的，单期效用 $u$ 是基数的，只在正线性变换下保持排序。
    \item 习惯以自己的过去消费为参照，攀比以他人消费为参照；参考点不对称、非线性概率权重等也可拓展标准模型。后文的组合选择与均衡定价主要使用 (6.3) 或 (6.4)。
    \item 圣彼得堡悖论区分期望金额与期望效用；CE 与包含初始财富和费用的最高入场费须分别计算。
    \item 阿莱斯例中的两对选择若反向排序，就与共同后果独立性矛盾；秩依效用、前景理论和心理账户分别强调概率权重、参考点及框架等作用，不能混为同一个假设。
\end{enumerate}
\end{tcolorbox}
\fi

% !TeX root = ../../金融经济学_王江_讲义及习题.tex
% 06_期望效用函数 课后习题（由章节文件拆分）

\begin{problem}
\problemno{6.1} 证明如下形式的期望效用函数：
\[
    U(c)=u_{0}(c_{0})+\sum_{\omega\in\Omega}\pi_{\omega}u_{1}(c_{1\omega})
\]
满足独立性公理。
\end{problem}
\begin{solution}
设两个计划 $c,c'$ 在状态 $\omega'$ 的消费路径同为 $(c_0,x)$，并且 $c\succeq c'$。由效用表示，
\[
 U(c)-U(c')=
 \sum_{\omega\ne\omega'}\pi_\omega
 \bigl[u_1(c_{1\omega})-u_1(c'_{1\omega})\bigr]\ge0.
\]
把共同路径换成 $(c_0,y)$，记新计划为 $d,d'$。两边共同的 $u_0(c_0)$ 和 $\pi_{\omega'}u_1(y)$ 仍抵消，故
\[
 U(d)-U(d')=U(c)-U(c')\ge0.
\]
因此 $d\succeq d'$，排序保持不变，独立性公理成立。
\end{solution}

\begin{problem}
\problemno{6.2*} 对于由练习 6.1 给出的效用函数表示的偏好，证明连续性意味着 $u_{0}(\cdot)$ 和 $u_{1}(\cdot)$ 都是连续的。
\end{problem}
\begin{solution}
由喜多厌少，$u_0,u_1$ 严格递增，且 $\sum_\omega\pi_\omega=1$。

考虑未来各状态消费均为 $y$ 的计划，其效用为
\[
 F(x,y)=u_0(x)+u_1(y).
\]
原偏好限制在这些计划上仍连续。下面证明：若 $F(x,y)=f(x)+g(y)$ 表示连续偏好，且 $f,g$ 严格递增，则 $f$ 连续。

单调函数在内点具有左右极限，且只有至多可数个不连续点。取 $g$ 的一个内点连续点 $y$。
\begin{enumerate}[label=(\roman*)]
\item 若 $f$ 在 $x$ 有左跳跃，记 $\delta=f(x)-f(x-)>0$。由 $g$ 在 $y$ 连续且严格递增，可取 $h>0$ 使
\[
 0<g(y)-g(y-h)<\delta.
\]
取 $x_n\uparrow x$，则
\[
 F(x_n,y)\le f(x-)+g(y)<f(x)+g(y-h)=F(x,y-h).
\]
由下等值集闭性，令 $n\to\infty$ 得 $(x,y)\preceq(x,y-h)$，与 $g$ 严格递增矛盾。
\item 若 $f$ 在 $x$ 有右跳跃，记 $\delta=f(x+)-f(x)>0$。取 $h>0$ 使
\[
 0<g(y+h)-g(y)<\delta.
\]
对 $x_n\downarrow x$，有 $F(x_n,y)>F(x,y+h)$。上等值集闭性推出 $(x,y)\succeq(x,y+h)$，同样矛盾。
\end{enumerate}
所以左右跳跃均不存在，$f$ 连续；在定义域端点只需相应的单侧论证。令 $(f,g)=(u_0,u_1)$ 得 $u_0$ 连续；交换二者，得 $u_1$ 连续。

\end{solution}

\begin{problem}
\problemno{6.3**} 如果 $a\succeq b$ 意味着对于 $\forall\,\alpha\in[0,1]$ 都有 $\alpha a+(1-\alpha)b\succeq b$，那么称偏好关系 $\succ$ 是弱凸的（weakly convex）。这与第 2 章中的强凸性有所不同。假设偏好关系是连续的并且有练习 6.1 给出的期望效用函数表示（因而 $u_{0}(\cdot)$ 和 $u_{1}(\cdot)$ 都是连续函数）。证明此时强凸性包含弱凸性。
\end{problem}
\begin{solution}
按本题的强凸性条件，若 $a\succ b$，则对任意 $0<\alpha<1$，有
\[
 \alpha a+(1-\alpha)b\succ b.
\]
要证弱凸性，只须补上 $a\sim b$ 的情形；$\alpha=0,1$ 以及 $a\succ b$ 的情形都已成立。

设 $a\sim b$，$0<\alpha<1$，记 $m=\alpha a+(1-\alpha)b$。反设 $m\prec b$，则
\[
 U(m)<U(a)=U(b).
\]
由于 $U$ 连续，沿着从 $a$ 到 $m$ 的线段，由介值定理可取
\[
 d=(1-t)a+tm,\quad 0<t<1,\qquad
 U(m)<U(d)<U(b).
\]
于是 $b\succ d$。另一方面，$m$ 位于 $d$ 与 $b$ 之间。具体地，令
\[
 \beta=\frac{\alpha}{1-t+t\alpha}\in(0,1),
 \qquad m=\beta d+(1-\beta)b.
\]
由 $b\succ d$ 和强凸性，应有 $m\succ d$，与 $U(m)<U(d)$ 矛盾。因此 $m\succeq b$，弱凸性成立。

\end{solution}

\begin{problem}
\problemno{6.4**} 假设定义在随机支付上的效用函数具有练习 6.1 中的形式，其中 $\Omega$ 是任意的状态空间而 $P$ 是相应的概率测度，并且偏好满足弱凸性。
\begin{enumerate}[label=(\alph*)]
    \item 证明 $u_{0}(\cdot)$ 和 $u_{1}(\cdot)$ 都是拟凹的。
    \item 它们是否必为凹的？若否，给出反例。
\end{enumerate}
\end{problem}
\begin{solution}
\begin{enumerate}[label=(\alph*)]
    \item $u$ 拟凹是指：若 $u(x)\ge u(y)$，则对任意 $\lambda\in[0,1]$ 有 $u(\lambda x+(1-\lambda)y)\ge u(y)$。也就是较高的那个函数值不会被凸组合拉到更低的值之下。只证 $u_{1}$，$u_{0}$ 同样。任取 $x,y$ 与 $\lambda\in[0,1]$，并取定 $c_{0}$ 以及 $1$ 期除某一状态 $\omega$ 外的消费。令 $a$ 在状态 $\omega$ 的消费为 $x$、$b$ 在状态 $\omega$ 的消费为 $y$，其余消费与 $a$ 相同。若 $u_{1}(x)\ge u_{1}(y)$，则 $U(a)\ge U(b)$，弱凸性给出 $U(\lambda a+(1-\lambda)b)\ge U(b)$，即 $u_{1}(\lambda x+(1-\lambda)y)\ge u_{1}(y)$。故 $u_{1}$ 拟凹。
    \item 不必凹。取 $u_{1}\equiv 0$、$u_{0}(c)=c^{3}$（$c\ge 0$），概率任意。此时 $U$ 只依赖于 $c_{0}$ 且对 $c_{0}$ 严格递增，故 $U(a)\ge U(b)$ 意味着 $a$ 的 $0$ 期消费不低于 $b$，其凸组合亦然，弱凸性成立。但 $u_{0}(c)=c^{3}$ 在 $[0,\infty)$ 上是凸的，不是凹的。
\end{enumerate}
\end{solution}


